\documentclass[12pt, a4paper]{article}
\usepackage[utf8]{inputenc}
\usepackage[T1]{fontenc}
\usepackage[french]{babel}
\usepackage[margin=2.5cm]{geometry}
\usepackage{amsmath, amssymb, amsthm, mathtools}
\usepackage{listings}
\usepackage{xcolor}
\usepackage{tcolorbox}
\tcbuselibrary{theorems, skins, breakable}
\usepackage{booktabs}
\usepackage{fancyhdr}
\usepackage{hyperref}
\usepackage{tikz}
\usetikzlibrary{arrows.meta, positioning, shapes.geometric}

\definecolor{suva_blue}{RGB}{30, 100, 180}
\definecolor{code_bg}{RGB}{245, 247, 250}
\definecolor{code_kw}{RGB}{0, 80, 200}
\definecolor{code_str}{RGB}{180, 40, 40}
\definecolor{code_cmt}{RGB}{80, 130, 80}
\definecolor{warn_bg}{RGB}{255, 248, 220}
\definecolor{success_bg}{RGB}{230, 255, 235}
\definecolor{falcon_color}{RGB}{180, 40, 120}
\definecolor{falcon_bg}{RGB}{255, 240, 250}

\lstdefinestyle{python}{
  language=Python,
  basicstyle=\ttfamily\small,
  keywordstyle=\color{code_kw}\bfseries,
  stringstyle=\color{code_str},
  commentstyle=\color{code_cmt}\itshape,
  backgroundcolor=\color{code_bg},
  frame=single, rulecolor=\color{gray!40},
  numbers=left, numberstyle=\tiny\color{gray},
  breaklines=true, tabsize=4, showstringspaces=false,
}
\lstdefinestyle{output}{
  basicstyle=\ttfamily\small,
  backgroundcolor=\color{gray!10},
  frame=single, rulecolor=\color{gray!50},
  breaklines=true,
}

\newtcolorbox{defbox}[2][]{colback=success_bg,colframe=green!50!black,
  fonttitle=\bfseries,title={Définition : #2},#1}
\newtcolorbox{falconbox}[2][]{colback=falcon_bg,colframe=falcon_color,
  fonttitle=\bfseries,title={Falcon : #2},#1}
\newtcolorbox{exercice}[2][]{colback=white,colframe=suva_blue!60,
  fonttitle=\bfseries,title={Exercice #2},#1,breakable}
\newtcolorbox{solution}[1][]{colback=gray!5,colframe=gray!40,
  fonttitle=\bfseries\itshape,title={Solution},#1,breakable}
\newtcolorbox{importantbox}[1][]{colback=suva_blue!10,colframe=suva_blue,
  fonttitle=\bfseries,title={Connexion avec SuVa},#1}
\newtcolorbox{warningbox}[1][]{colback=warn_bg,colframe=orange!80!black,
  fonttitle=\bfseries,title={Attention},#1}
\newtcolorbox{compbox}[1][]{colback=gray!5,colframe=gray!60,
  fonttitle=\bfseries,title={Comparaison Falcon vs Dilithium},#1}

\pagestyle{fancy}
\fancyhf{}
\fancyhead[L]{\color{suva_blue}\textbf{SuVa --- Chapitre 7}}
\fancyhead[R]{\color{gray}Falcon / FN-DSA}
\fancyfoot[L]{\footnotesize\color{gray}Mohamed Lamine MBENGUE}
\fancyfoot[C]{\thepage}
\fancyfoot[R]{\footnotesize\color{gray}portfolio-alamine.web.app}
\renewcommand{\headrulewidth}{0.4pt}
\renewcommand{\footrulewidth}{0.3pt}

% ══════════════════════════════════════════════════════════════════════
\begin{document}

\begin{center}
  {\color{suva_blue}\rule{\linewidth}{2pt}}\\[0.5cm]
  {\Huge\bfseries\color{suva_blue} Chapitre 7}\\[0.3cm]
  {\large\bfseries Falcon / FN-DSA (NIST FIPS 206)}\\[0.2cm]
  {\small\color{gray} La prochaine étape du projet SuVa --- la cible production}\\[0.3cm]
  {\color{suva_blue}\rule{\linewidth}{2pt}}
\end{center}

\vspace{0.3cm}
\begin{tcolorbox}[colback=suva_blue!5, colframe=suva_blue!40,
  left=6pt, right=6pt, top=4pt, bottom=4pt]
\small
\begin{tabular}{ll}
  \textbf{Auteur}      & Mohamed Lamine MBENGUE \\
  \textbf{Spécialités} & Sécurité · DevOps · Dev Full Stack Web \& Mobile \\
  \textbf{Contact}     & +221 78 178 44 36 · mbenguemohamedlamine2022@gmail.com \\
  \textbf{Portfolio}   & \href{https://portfolio-alamine.web.app}{\texttt{portfolio-alamine.web.app}} \\
\end{tabular}
\end{tcolorbox}
\vspace{0.3cm}

\tableofcontents
\newpage

%══════════════════════════════════════════════════════════════════════
\section{Pourquoi Falcon est la cible de SuVa}
%══════════════════════════════════════════════════════════════════════

\begin{center}
\begin{tabular}{lrrrll}
\toprule
\textbf{Schéma} & \textbf{Gas} & \textbf{Sig (bytes)} & \textbf{pk (bytes)}
  & \textbf{NIST} & \textbf{Sécurité} \\
\midrule
ECDSA (actuel) & 3K & 65 & 33 & — & Cassable quantique \\
\textbf{Falcon-512} & \textbf{350K} & \textbf{666} & \textbf{897} & FIPS 206 & Lattice (NTRU) \\
Dilithium2 & 13.5M & 2420 & 1312 & FIPS 204 & Lattice (MLWE) \\
SPHINCS+ & 50--500M & 7856+ & 32 & FIPS 205 & Hash-based \\
\bottomrule
\end{tabular}
\end{center}

\textbf{Falcon est $\sim 20\times$ moins cher que Dilithium et $\sim 116\times$
plus compact en signature.} C'est pourquoi SuVa vise Falcon comme schéma
principal pour les transactions quotidiennes.

%══════════════════════════════════════════════════════════════════════
\section{NTRU --- La base mathématique de Falcon}
%══════════════════════════════════════════════════════════════════════

\subsection{Différence fondamentale avec Dilithium}

\begin{compbox}
\begin{center}
\begin{tabular}{lll}
\toprule
& \textbf{Dilithium} & \textbf{Falcon} \\
\midrule
Anneau & $R_q = \mathbb{Z}_q[x]/(x^n+1)$ & $R_q = \mathbb{Z}_q[x]/(x^n+1)$ \\
Problème dur & MLWE, MSIS & NTRU, SIS \\
Clé publique & $(rho, t_1)$ & $h = g \cdot f^{-1}$ \\
Structure clé secrète & $(s_1, s_2)$ petits vecteurs & $(f, g)$ polynômes courts \\
Keygen & Rapide & Lent (FFT Gram-Schmidt) \\
Signe & Boucle de rejet & Gaussian sampling \\
Signature & $(c, z, h)$ & $(s_1, s_2)$ court \\
\bottomrule
\end{tabular}
\end{center}
\end{compbox}

\subsection{Le problème NTRU}

\begin{defbox}{Problème NTRU}
Dans $R_q = \mathbb{Z}_q[x]/(x^n + 1)$, connaissant~:
$$h = g \cdot f^{-1} \pmod{q}$$
où $f, g$ ont des coefficients petits (Gaussiens), retrouver $f$ ou $g$.

La clé publique est $h$, la clé secrète est $(f, g)$.
\end{defbox}

\begin{importantbox}
\textbf{Paramètres Falcon-512 (niveau NIST I)~:}
\begin{itemize}
  \item $n = 512$, $q = 12\,289$
  \item $f, g$ : coefficients Gaussiens d'écart-type $\sigma \approx 1.17\sqrt{q/2n}$
  \item Sécurité : $\approx 103$ bits (classique)
\end{itemize}

\textbf{Paramètres Falcon-1024 (niveau NIST V)~:}
\begin{itemize}
  \item $n = 1024$, $q = 12\,289$
  \item Sécurité : $\approx 207$ bits (classique)
\end{itemize}
\end{importantbox}

%══════════════════════════════════════════════════════════════════════
\section{KeyGen Falcon --- Génération de la trappe NTRU}
%══════════════════════════════════════════════════════════════════════

\subsection{Vue d'ensemble}

\begin{lstlisting}[style=python]
# KeyGen Falcon (pseudo-code)
def falcon_keygen(n=512, q=12289):
    # 1. Generer f, g avec des coefficients courts (Gaussiens)
    while True:
        f = sample_gaussian_poly(n, sigma=1.17 * (q/(2*n))**0.5)
        g = sample_gaussian_poly(n, sigma=1.17 * (q/(2*n))**0.5)

        # f doit etre inversible dans R_q
        if not is_invertible(f, q, n):
            continue

        # 2. Calculer h = g * f^(-1) mod q
        f_inv = poly_inverse(f, q, n)
        h = poly_mul(g, f_inv, q, n)

        # 3. Generer la trappe NTRU (pour le signing)
        # Calculer (F, G) tels que f*G - g*F = q (equation NTRU)
        F, G = solve_ntru(f, g, q, n)
        # C'est la partie difficile : utilise FFT sur les entiers

        if F is not None and G is not None:
            break

    pk = h          # cle publique : h (897 bytes pour Falcon-512)
    sk = (f, g, F, G)  # cle secrete : la trappe NTRU
    return pk, sk
\end{lstlisting}

\begin{warningbox}
La génération de clés Falcon est \textbf{significativement plus complexe}
que Dilithium car elle nécessite la décomposition LDL (via FFT sur les
entiers de Gauss). C'est pourquoi on ne génère les clés qu'une fois
et on les réutilise.
\end{warningbox}

%══════════════════════════════════════════════════════════════════════
\section{Sign Falcon --- Gaussian Sampling}
%══════════════════════════════════════════════════════════════════════

Falcon utilise le \textbf{Gaussian sampling} (pas de boucle de rejet
comme Dilithium) pour produire des signatures très courtes.

\subsection{Algorithme de signature}

\begin{lstlisting}[style=python]
# Sign Falcon (pseudo-code tres simplifie)
def falcon_sign(sk, message, n=512, q=12289):
    f, g, F, G = sk

    # 1. Hacher le message -> point cible
    salt = random_bytes(40)
    c = hash_to_point(message, salt, n, q)
    # c est un polynome dans R_q

    # 2. Gaussian preimage sampling
    # Trouver (s1, s2) courts tels que s1*h + s2 = c mod q
    # avec (s1, s2) distribues gaussianement
    # On utilise la trappe (f, g, F, G) pour echantillonner

    # La trappe permet de centrer la Gaussienne autour de la
    # "bonne" solution (la plus courte)
    B = [[g, -f], [G, -F]]  # base NTRU
    # Fast Fourier Sampling dans le reseau NTRU
    s1, s2 = ffsampling_fft(c, B, n, q)

    # Verifier la norme
    sig_norm = l2_norm(s1) + l2_norm(s2)
    # La norme doit etre <= 1.1 * beta_NTRU
    # Si trop grande, recommencer (rare)

    # Signature = (salt, s2) -- s1 peut etre retrouve de h et s2
    return bytes(salt) + pack_int12(s2)
\end{lstlisting}

\begin{importantbox}
\textbf{Gaussian Sampling vs Rejection Sampling~:}
\begin{itemize}
  \item Falcon~: la signature suit une distribution Gaussienne autour
        de la solution optimale. Presque toujours valide du premier coup.
  \item Dilithium~: on échantillonne uniformément et on rejette si trop grand.
        Nécessite en moyenne 4.25 itérations.
\end{itemize}
\textbf{Résultat~:} Falcon produit des signatures \textbf{3--4x plus courtes}
(666 vs 2420 bytes) et s'exécute en \textbf{1 itération}.
\end{importantbox}

%══════════════════════════════════════════════════════════════════════
\section{Verify Falcon --- Vérification simple}
%══════════════════════════════════════════════════════════════════════

La vérification Falcon est \textbf{beaucoup plus simple} que la signature~:

\begin{lstlisting}[style=python]
# Verify Falcon
def falcon_verify(pk, message, signature, n=512, q=12289):
    h = pk  # cle publique

    # Decoder la signature
    salt = signature[:40]
    s2 = unpack_int12(signature[40:])

    # Recomputer s1 : s1 = h^(-1) * (c - s2) mod q
    c = hash_to_point(message, salt, n, q)
    s1 = poly_mul(poly_inverse(h, q, n), c - s2, q, n)

    # Verifier la norme : ||(s1, s2)||_2 <= beta
    norm_sq = sum(x*x for x in s1) + sum(x*x for x in s2)
    # beta^2 est fixe dans les parametres
    return norm_sq <= BETA_SQ
\end{lstlisting}

\textbf{C'est la raison du faible coût gas~:} la vérification ne fait
qu'une multiplication polynomiale + vérification de norme.

%══════════════════════════════════════════════════════════════════════
\section{ETHFALCON --- implémentation EVM disponible}
%══════════════════════════════════════════════════════════════════════

\begin{importantbox}
ZKNox a déjà implémenté Falcon pour l'EVM (ETHFALCON, disponible sur
GitHub : \texttt{ZKNoxHQ/ETHFALCON}).

Pour SuVa, la prochaine étape est~:
\begin{enumerate}
  \item Étudier ETHFALCON (Solidity + Python)
  \item L'intégrer dans le smart wallet SuVa (ERC-4337)
  \item Adapter le SDK Python pour signer avec Falcon
  \item Générer les vecteurs de test Python → Solidity
\end{enumerate}
\end{importantbox}

\subsection{Structure anticipée de l'intégration Falcon dans SuVa}

\begin{lstlisting}[style=python]
# Futur SDK SuVa (a construire)

# Cote Python (off-chain)
from suva_sdk import FalconWallet

wallet = FalconWallet()
pk, sk = wallet.generate_keys()     # Falcon-512 keygen
print(f"Cle publique : {len(pk)} bytes")  # 897 bytes

# Signer une transaction
tx = {"to": "0x...", "value": 1000, "nonce": 42}
sig = wallet.sign_transaction(sk, tx)  # Falcon sign
print(f"Signature : {len(sig)} bytes")   # ~666 bytes

# Cote Solidity (on-chain) - contracte SuVa Wallet
# validateUserOp() appelle verify() du contrat ETHFALCON
# ~350K gas
\end{lstlisting}

%══════════════════════════════════════════════════════════════════════
\section{Roadmap Falcon dans SuVa}
%══════════════════════════════════════════════════════════════════════

\begin{center}
\begin{tikzpicture}[
  step/.style={draw=suva_blue, rounded corners, fill=blue!5,
    minimum width=5cm, minimum height=1cm, text width=4.8cm,
    align=left, font=\small, inner sep=6pt},
  done/.style={step, fill=success_bg, draw=green!50!black},
  todo/.style={step, fill=warn_bg, draw=orange!80},
  critical/.style={step, fill=red!10, draw=red!60},
  arrow/.style={->, thick}
]
  \node[done] (a) {
    \textbf{[FAIT]} Dilithium Python\\
    keygen + sign + verify
  };
  \node[done, below=0.6cm of a] (b) {
    \textbf{[FAIT]} Dilithium Solidity (ZKNox)\\
    vérification on-chain
  };
  \node[todo, below=0.6cm of b] (c) {
    \textbf{[TODO]} Intégrer ETHFALCON\\
    dans le smart wallet SuVa
  };
  \node[todo, below=0.6cm of c] (d) {
    \textbf{[TODO]} SDK client Falcon\\
    (Python, puis JS/TypeScript)
  };
  \node[critical, below=0.6cm of d] (e) {
    \textbf{[TODO]} Vault contracts\\
    wrapper USDT/USDC → suUSDT
  };
  \node[critical, below=0.6cm of e] (f) {
    \textbf{[TODO]} Audit de sécurité\\
    avant tout déploiement
  };

  \draw[arrow] (a) -- (b);
  \draw[arrow] (b) -- (c);
  \draw[arrow] (c) -- (d);
  \draw[arrow] (d) -- (e);
  \draw[arrow] (e) -- (f);
\end{tikzpicture}
\end{center}

%══════════════════════════════════════════════════════════════════════
\section{Comparaison finale : Dilithium vs Falcon vs SPHINCS+}
%══════════════════════════════════════════════════════════════════════

\begin{center}
\begin{tabular}{llll}
\toprule
& \textbf{Dilithium/ML-DSA} & \textbf{Falcon/FN-DSA} & \textbf{SPHINCS+} \\
\midrule
\textbf{Sécurité} & MLWE + MSIS & NTRU + SIS & Uniquement hash \\
\textbf{Hypothèse} & Module-LWE & NTRU & Fonctions de hachage \\
\textbf{Quantum safe} & Oui & Oui & Oui (le plus sûr) \\
\textbf{Pk size (L1)} & 1312 bytes & 897 bytes & 32 bytes \\
\textbf{Sig size (L1)} & 2420 bytes & 666 bytes & 7856 bytes \\
\textbf{Keygen speed} & Rapide & Lent & Rapide \\
\textbf{Sign speed} & Moyen & Rapide (1 iter) & Lent \\
\textbf{Verify speed} & Moyen & Rapide & Lent \\
\textbf{Gas EVM} & 13.5M & \textbf{350K} & 50--500M \\
\textbf{Rôle SuVa} & Recherche & \textbf{Production} & Fallback \\
\bottomrule
\end{tabular}
\end{center}

%══════════════════════════════════════════════════════════════════════
\section{Exercices}
%══════════════════════════════════════════════════════════════════════

\begin{exercice}{1 --- Vérifier les tailles de Falcon}
Falcon-512 produit des signatures de $\lceil 666 \rceil$ bytes.
\begin{enumerate}[label=\alph*)]
  \item La signature contient un sel de 40 bytes + $s_2$ encodé sur
        $\lceil n \log_2(2 \cdot 6146 + 1) \rceil$ bits.
        Pour $n = 512$, calcule la taille théorique de $s_2$.
  \item Pourquoi la signature Falcon est-elle $3.6\times$ plus petite
        que Dilithium2 malgré que $n$ soit $2\times$ plus grand ?
\end{enumerate}
\end{exercice}

\begin{solution}
\begin{lstlisting}[style=python]
import math

# a) Taille de s2
n = 512
# Les coefficients de s2 sont dans [-6146, 6146] approximativement
# encoding compact sur 12 bits par coeff (~12*512/8 = 768 bytes)
# + sel 40 bytes + compression
bits_s2 = n * 12  # approximation
bytes_s2 = bits_s2 // 8
print(f"s2 approximatif : {bytes_s2} bytes")  # ~768
print(f"+ sel 40 bytes = {bytes_s2 + 40} bytes")

# b) Falcon est plus compact car :
# - Distribution Gaussienne -> coefficients plus petits en moyenne
# - Pas de z (qui necessite 18 bits/coeff dans Dilithium)
# - Pas de hint h
# - Pas de c_tilde 32 bytes (le sel remplace)
# Dilithium : c(32) + z(4*256*18/8=2304) + h(~84) = 2420 bytes
# Falcon-512 : sel(40) + s2(~626) = ~666 bytes
\end{lstlisting}
\end{solution}

\begin{exercice}{2 --- Explorer ETHFALCON}
ETHFALCON est disponible sur \texttt{github.com/ZKNoxHQ/ETHFALCON}.
\begin{enumerate}[label=\alph*)]
  \item Clone le repo et explore sa structure.
  \item Quels sont les contrats Solidity clés ?
  \item Y a-t-il une implémentation Python de référence ?
  \item Comment les tests sont-ils organisés ?
  \item Quel est le gas cost mesuré par les tests ?
\end{enumerate}
\end{exercice}

\begin{exercice}{3 --- Plan d'intégration}
Conçois le plan d'intégration de Falcon dans SuVa~:
\begin{enumerate}[label=\alph*)]
  \item Quels fichiers Python du projet actuel devront être adaptés
        pour Falcon ?
  \item Quels contrats Solidity de ETHDILITHIUM sont réutilisables pour
        Falcon (NTT, utils...) ?
  \item Comment les test vectors Python → Solidity seront-ils générés
        pour Falcon ?
  \item Quelles modifications faudra-t-il apporter au smart wallet
        ERC-4337 ?
\end{enumerate}
\end{exercice}

\begin{exercice}{4 --- SPHINCS+ comme fallback}
SPHINCS+ est le fallback "nuclear bunker" de SuVa.
Sa sécurité repose uniquement sur les fonctions de hachage.
\begin{enumerate}[label=\alph*)]
  \item Pourquoi SPHINCS+ est-il le plus sûr des trois schémas ?
  \item Pourquoi son coût gas est-il si élevé (50--500M) ?
  \item Dans quel scénario SuVa utiliserait-il SPHINCS+ plutôt que Falcon ?
  \item SPHINCS+ a des variantes "small" (petite signature, lent) et
        "fast" (grande signature, rapide). Laquelle préférer pour SuVa ?
\end{enumerate}
\end{exercice}

\end{document}
